\documentclass[11pt,a4paper]{article}
\usepackage[T1]{fontenc}
\usepackage{lmodern,amsmath,amssymb,amsthm,microtype}
\usepackage[margin=27mm]{geometry}
\usepackage[colorlinks=true,linkcolor=blue,citecolor=blue,urlcolor=blue]{hyperref}
\setlength{\emergencystretch}{2em}
\newtheorem{theorem}{Theorem}
\newtheorem{lemma}[theorem]{Lemma}
\newcommand{\Q}{\mathbb Q}
\begin{document}
\title{A prepared input with a nontrivial one-coordinate covering obstruction}
\author{}\date{17 November 2026}\maketitle

\begin{theorem}\label{thm:main}
The rational number $d=112/225$ satisfies
\[
 0<d<1,\quad\operatorname{den}(d)=15^2,\quad
 d=7(4/15)^2,\quad v_2(d)=4.
\]
The smooth genus-one curve
\[
 C_d:\quad y^2=d\bigl((1-d)^2-s^4\bigr)
\]
has a point over every completion of $\Q$, but no rational point. Consequently the specialization $q_0=0$ of the intrinsic two-coordinate construction does not supply witnesses for every prepared input.
\end{theorem}

\section{The covering and its local points}
The changes $t=15s$ and $z=(3375/4)y$ identify $C_d$ with
\begin{equation}\label{eq:quartic}
 z^2=7(113^2-t^4).
\end{equation}
The Jacobian is $E'_{791}:Y^2=X^3+4\cdot791^2X$.
The map
\begin{equation}\label{eq:cover}
 (t,z)\longmapsto(X,Y)=(-7t^2,7tz)
\end{equation}
has degree two onto
\[
 E=E_{791}:\quad Y^2=X^3-791^2X.
\]
A rational point of \eqref{eq:quartic} would therefore give a point of $E(\Q)$ with abscissa class $[-7]$. There are no rational points of \eqref{eq:quartic} with $t=0$, since $7$ is nonsquare; no rational points with $z=0$, since $113$ is nonsquare; and no rational points at infinity, since $-7$ is nonsquare.

The local solubility can be checked directly. At $2$, $-7$ is a square, and large $2$-adic $t$ gives an affine point by Hensel's lemma. At $7$, $113\equiv1\pmod7$ is a local square, so $t^2=113,z=0$ gives a point. At $113$, $7$ is a local square since $32^2\equiv7\pmod{113}$; taking $t=0$ gives a point. Outside $2,7,113$ the quartic has good genus-one reduction, and the Hasse bound followed by Hensel lifting gives a local point. The real place is immediate. At every completion a smooth local neighbourhood avoids the finite set $tz=0$ and infinity if desired.

\section{The Selmer group and the Cassels pairing convention}
We use the standard full $2$-descent convention of Wang \cite[Sections 2.1--2.2]{W} and Zhang \cite[Sections 2.1--2.2]{Z}. For $n>0$ odd and squarefree, an element $\Lambda=(d_1,d_2,d_3)$ of $\operatorname{Sel}_2(E_n)$ is represented by the intersection
\begin{align*}
 H_1&=-nt^2+d_2u_2^2-d_3u_3^2=0,\\
 H_2&=-nt^2+d_3u_3^2-d_1u_1^2=0,\\
 H_3&=2nt^2+d_1u_1^2-d_2u_2^2=0.
\end{align*}
Only two equations are independent. A point $(X,Y)$ maps to the square-class triple $(X-n,X+n,X)$. The images of $O,(n,0),(-n,0),(0,0)$ are respectively
\begin{equation}\label{eq:torsion}
 (1,1,1),\quad (2,2n,n),\quad(-2n,2,-n),\quad(-n,n,-1).
\end{equation}
For each $i$, choose a rational point on the conic $H_i=0$ and a tangent linear form $L_i$. Then the multiplicative Cassels pairing is
\begin{equation}\label{eq:pairing}
 \langle\Lambda,\Lambda'\rangle
 =\prod_v\prod_{i=1}^3 (L_i(P_v),d_i')_v,
 \qquad P_v\in D_\Lambda(\Q_v),
\end{equation}
where $(\ ,\ )_v$ is the quadratic Hilbert symbol. It is alternating on the pure Selmer group $\operatorname{Sel}_2(E_n)/E_n(\Q)[2]$. Its radical contains the image of $E_n(\Q)/2E_n(\Q)$ modulo the torsion images. This follows, for example, from the stronger statement that the radical is the image of the $4$-Selmer group under multiplication by two, modulo the torsion images \cite[Section 2.2]{W}.

For $n=791=7\cdot113$, order the primes as $(7,113)$. The off-diagonal Legendre symbols are both $1$, so the Monsky matrices in Wang's convention are
\[
 A=0,\qquad D_2=0,\qquad D_{-2}=\begin{pmatrix}1&0\\0&0\end{pmatrix},
 \qquad
 M=\begin{pmatrix}1&0&0&0\\0&0&0&0\\0&0&0&0\\0&0&0&0\end{pmatrix}.
\]
The pure $2$-Selmer group therefore has dimension three, with basis
\begin{equation}\label{eq:basis}
 (113,1,113),\qquad (1,7,7),\qquad(1,113,113).
\end{equation}

\section{An exact nonzero pairing}
Take
\[
 \Lambda=(1,7,7),\qquad \Lambda'=(113,1,113).
\]
The following table lists rational points on the three conics and primitive integral tangent forms. The omitted coordinate $u_i$ is irrelevant to $H_i$.
\[
\begin{array}{c|c|l}
 &\text{rational conic point}&L_i\\\hline
 H_1&(t,u_2,u_3)=(0,1,1)&u_2-u_3\\
 H_2&(t,u_1,u_3)=(1,28,15)&-113t-4u_1+15u_3\\
 H_3&(t,u_1,u_2)=(1,21,17)&226t+3u_1-17u_2
\end{array}
\]
The original equations here are
\[
 H_1=-791t^2+7u_2^2-7u_3^2,
 \quad H_2=-791t^2+7u_3^2-u_1^2,
 \quad H_3=1582t^2+u_1^2-7u_2^2.
\]
Because the middle coordinate of $\Lambda'$ is $1$, formula \eqref{eq:pairing} reduces to
\begin{equation}\label{eq:two}
 \langle\Lambda,\Lambda'\rangle
 =\prod_v\bigl((u_2-u_3)(226t+3u_1-17u_2),113\bigr)_v.
\end{equation}

At $v=\infty$ this symbol is $1$, since $113>0$. At $2$ it is $1$ because $113\equiv1\pmod8$; at $7$ it is $1$ because $113\equiv1\pmod7$.

Here is an explicit all-prime verification of vanishing outside $2,7,113$, including the primes $3$ and $17$ appearing in tangent coefficients. At an odd prime $p\notin\{7,113\}$, scale $P_p$ to have integral primitive projective coordinates. On $D_\Lambda$ we have
\begin{align}
 (u_2-u_3)(u_2+u_3)&=113t^2,\label{eq:identity1}\\
 (226t+3u_1-17u_2)(226t-3u_1+17u_2)
 &=-\frac{(17u_1-21u_2)^2}{7}.\label{eq:identity3}
\end{align}
If both factors on the left of \eqref{eq:identity1} were divisible by $p$, then $u_2,u_3,t$, and hence $u_1$, would all be divisible by $p$. This contradicts primitivity. Thus any positive valuation of $L_1$ is even.

If both factors on the left of \eqref{eq:identity3} were divisible by $p$, their sum $452t$ would force $p\mid t$. The identity and $L_3\equiv0$ would then give
\[
 17u_1-21u_2\equiv0,\qquad3u_1-17u_2\equiv0\pmod p.
\]
The determinant of these equations is $-226=-2\cdot113$, a unit at $p$. Therefore $u_1,u_2$, and then $u_3$, would all be divisible by $p$, again impossible. It follows that any positive valuation of $L_3$ is even. Thus $v_p(L_1L_3)$ is even; since $113$ is a unit, the Hilbert symbol is $1$. Local points may always be moved away from the finitely many tangent zeros before evaluation. This argument requires no exclusions at $3$, $5$, or $17$.

At $p=113$, let $r\in\Q_{113}$ be the square root of $7$ with $r\equiv32\pmod{113}$. Take
\[
 P_{113}=(t,u_1,u_2,u_3)=(0,-r,-1,1).
\]
This satisfies all three conics exactly. Its tangent product is
\[
 L_1(P_{113})L_3(P_{113})=-2(17-3r)
 \equiv45\pmod{113}.
\]
It is a unit and
\[
 \left(\frac{45}{113}\right)=\left(\frac5{113}\right)=-1.
\]
Therefore the local symbol at $113$ is $-1$, and
\begin{equation}\label{eq:nonzero}
 \boxed{\langle\Lambda,\Lambda'\rangle=-1.}
\end{equation}

\section{Rank and abscissa classes}
A nonzero alternating pairing on a three-dimensional vector space has rank two and a one-dimensional radical. Since the Mordell--Weil image modulo torsion lies in that radical, \eqref{eq:nonzero} proves
\[
 \operatorname{rank}E(\Q)\le1.
\]
The exact point
\[
 P=(-784,-2940)\in E(\Q)
\]
has full descent triple
\[
 [-1575,7,-784]=[-7,7,-1].
\]
Multiplying by the torsion triple $(-791,791,-1)$ gives
\[
 (113,113,1),
\]
which is nonzero in the pure Selmer group \eqref{eq:basis}. In particular $P$ has nonzero image modulo $2E(\Q)$ and torsion. To justify the torsion qualification without assuming a torsion table, none of the three nonidentity triples in \eqref{eq:torsion} is trivial; hence no nonzero rational point of order two is divisible by two. Thus there is no rational $4$-torsion. All odd torsion is divisible by two, and the image of the torsion subgroup in full $2$-descent is exactly \eqref{eq:torsion}. This proves $P$ is non-torsion, and the rank is exactly one. It also proves that $P$ is an odd multiple of a generator modulo torsion.

The abscissa square-class homomorphism
\[
 \delta:E(\Q)\to\Q^\times/\Q^{\times2}
\]
sends $(0,0)$ to $[-1]$, $(791,0)$ to $[791]$, and $P$ to $[-1]$. It kills doubles. Since $P$ generates the free part modulo two, the full image is precisely
\[
 \delta(E(\Q))=\{[1],[-1],[791],[-791]\}.
\]
The class $[-7]$ is not in this set. The map \eqref{eq:cover} now excludes every rational point of \eqref{eq:quartic}, proving Theorem~\ref{thm:main}.

\section{Scope}
This counterexample concerns a prepared input, and not an asserted actual detector tuple. It rules out uniform completeness of the $q_0=0$ one-coordinate replacement on all prepared inputs. It does not rule out a restricted actual-trace family, a different one-coordinate relation, or individual rational witnesses for the original two-coordinate intrinsic block. It also demonstrates why local solubility and positive rank of the congruent curve do not suffice to force the specified abscissa class.

\begin{thebibliography}{9}
\bibitem{W} Z. Wang, \emph{Congruent elliptic curves with non-trivial Shafarevich--Tate groups}, Science China Mathematics 59 (2016), 2145--2166; arXiv:1511.03810, Sections 2.1--2.2, especially (2.1)--(2.3) and the radical statement after (2.3). \url{https://arxiv.org/pdf/1511.03810}
\bibitem{Z} S. Zhang, \emph{On non-congruent numbers as multiples of non-congruent numbers}, 28 August 2025 manuscript, Sections 2.1--2.2, formulas (2.1), (2.4)--(2.5), and the tangent-plane pairing. \url{https://zhangshenxing.github.io/publications/Zhang2025.pdf}
\end{thebibliography}
\end{document}
